% ch1_answers.tex -- Quick check answers for Module 5, Chapter 1 (Projectile Motion).
% Convention in every answer: origin at the launch point, up positive, a_y = -9.8 m s^-2.
\begin{chapteranswers}{1}{Projectile Motion}
  % ---- 1.1 The projectile model
  \item (a) $u_x = 15\cos\ang{60} = \qty{7.5}{m.s^{-1}}$ and
        $u_y = 15\sin\ang{60} = \qty{12.99}{m.s^{-1}}$:
        \final{$u_x = \qty{7.5}{m.s^{-1}}$, $u_y = \qty{13}{m.s^{-1}}$}.
        (b) The angle is measured from the \emph{vertical}, so sin and cos swap. The dart goes
        upwards, so $u_y$ is positive:
        $u_x = 8.0\sin\ang{20} = \qty{2.736}{m.s^{-1}}$ and
        $u_y = 8.0\cos\ang{20} = \qty{7.518}{m.s^{-1}}$:
        \final{$u_x = \qty{2.7}{m.s^{-1}}$, $u_y = \qty{7.5}{m.s^{-1}}$}.
        (Check: \ang{70} to the horizontal gives $8.0\cos\ang{70} = \qty{2.7}{m.s^{-1}}$.)
  \item Vertical motion first: $u_y = 0$ and $s_y = \qty{-20}{m}$. From
        $s_y = u_y t + \frac{1}{2}a_y t^2$: $-20 = \frac{1}{2}(-9.8)t^2$, so $-20 = -4.9t^2$.
        Divide both sides by $-4.9$: $t^2 = \dfrac{-20}{-4.9} = \qty{4.082}{s^2}$, and
        $t = \qty{2.020}{s}$ (\final{$t = \qty{2.0}{s}$}). Horizontal motion, same time:
        $s_x = u_x t = 6.0 \times 2.020 = \qty{12.12}{m}$
        (\final{$s_x = \qty{12}{m}$} from the base of the cliff).
  \item They land \textbf{at the same time}. Both start with $u_y = 0$, and both have the
        same vertical acceleration, \qty{9.8}{m.s^{-2}} down (air resistance ignored). The
        horizontal motion of B does not affect its vertical motion. So at every moment both
        balls have fallen the same distance.
  \item (a) The time \textbf{doubles}. For a horizontal launch $t = \sqrt{2h/g}$. Four times
        the height gives $\sqrt{4} = 2$ times the time. (b) The range \textbf{doubles} too.
        $u_x$ is the same and $R = u_x t$, so twice the time gives twice the range.
  % ---- 1.2 Height, time of flight, range and final velocity
  \item $u_y = 25\sin\ang{50} = \qty{19.15}{m.s^{-1}}$ and $u_x = 25\cos\ang{50} = \qty{16.07}{m.s^{-1}}$.
        $h_\mathrm{max} = \dfrac{u_y^2}{2g} = \dfrac{19.15^2}{2 \times 9.8} = \dfrac{366.7}{19.6}
        = \qty{18.71}{m}$ (\final{\qty{19}{m}}).
        $t_\mathrm{flight} = \dfrac{2u_y}{g} = \dfrac{2 \times 19.15}{9.8} = \qty{3.908}{s}$
        (\final{\qty{3.9}{s}}).
        At the top $v_y = 0$, so the speed is $v = v_x = u_x = \qty{16.07}{m.s^{-1}}$
        (\final{\qty{16}{m.s^{-1}}}).
  \item First change the units: $u = 64.8 \div 3.6 = \qty{18.0}{m.s^{-1}}$.
        Then use $R = \dfrac{u^2\sin 2\theta}{g}$ with $u^2 = 18.0^2 = 324$.
        At \ang{45}: $R = \dfrac{324 \times \sin\ang{90}}{9.8} = \qty{33.06}{m}$ (\final{\qty{33}{m}}).
        At \ang{30}: $R = \dfrac{324 \times \sin\ang{60}}{9.8} = \qty{28.63}{m}$ (\final{\qty{29}{m}}).
        At \ang{60}: $R = \dfrac{324 \times \sin\ang{120}}{9.8} = \qty{28.63}{m}$
        (\final{\qty{29}{m}}). \ang{30} and \ang{60} give the same range because they are
        complementary angles.
  \item The water lands at the same height as the nozzle, so use the level-ground formula
        $R = \dfrac{u^2\sin 2\theta}{g}$, rearranged:
        $u = \sqrt{\dfrac{Rg}{\sin 2\theta}} = \sqrt{\dfrac{6.5 \times 9.8}{\sin\ang{80}}}
        = \sqrt{\dfrac{63.7}{0.9848}} = \sqrt{64.68} = \qty{8.04}{m.s^{-1}}$.
        Answer: \final{$u = \qty{8.0}{m.s^{-1}}$}. (Check:
        $8.04^2 \times \sin\ang{80}/9.8 = \qty{6.5}{m}$.)
  \item $u_x = 12\cos\ang{25} = \qty{10.88}{m.s^{-1}}$, $u_y = 12\sin\ang{25} = \qty{5.071}{m.s^{-1}}$,
        and $s_y = \qty{-18}{m}$. Then $-18 = 5.071t + \frac{1}{2}(-9.8)t^2 = 5.071t - 4.9t^2$. Add 18 to both sides and
        multiply every term by $-1$: $4.9t^2 - 5.071t - 18 = 0$
        ($A = 4.9$, $B = -5.071$, $C = -18$).
        $B^2 - 4AC = (-5.071)^2 - 4(4.9)(-18) = 25.72 + 352.8 = 378.5$, and $\sqrt{378.5} = 19.46$.
        $t = \dfrac{5.071 \pm 19.46}{9.8}$. The positive root is $t = \dfrac{24.53}{9.8} = \qty{2.503}{s}$
        (the other root, \qty{-1.47}{s}, is before the launch): \final{$t = \qty{2.5}{s}$}.
        Range: $R = u_x t = 10.88 \times 2.503 = \qty{27.2}{m}$, so \final{$R = \qty{27}{m}$}.
        Final velocity: $v_x = \qty{10.88}{m.s^{-1}}$;
        $v_y^2 = u_y^2 + 2a_y s_y = 25.72 + 2(-9.8)(-18) = 25.72 + 352.8 = \qty{378.5}{m^2.s^{-2}}$,
        so $v_y = \qty{-19.46}{m.s^{-1}}$ (down). $v = \sqrt{10.88^2 + 19.46^2} = \sqrt{118.4 + 378.5}
        = \qty{22.29}{m.s^{-1}}$ and $\alpha = \tan^{-1}(19.46 \div 10.88) = \ang{60.8}$.
        Answer: \final{$v = \qty{22}{m.s^{-1}}$ at \ang{61} below the horizontal}.
        (Energy check: $\sqrt{12^2 + 2 \times 9.8 \times 18} = \qty{22.29}{m.s^{-1}}$ \cmark.)
  \item First the time at the posts, then the height at that time. Keep extra digits:
        $u_x = 20\cos\ang{35} = \qty{16.383}{m.s^{-1}}$ and $u_y = 20\sin\ang{35} = \qty{11.4715}{m.s^{-1}}$.
        Time at the posts: $t = \dfrac{s_x}{u_x} = \dfrac{30}{16.383} = \qty{1.8312}{s}$.
        Height then: $s_y = u_y t + \frac{1}{2}a_y t^2 = 11.4715 \times 1.8312 + \frac{1}{2}(-9.8)(1.8312)^2$,
        so $s_y = 21.01 - 16.43 = \qty{4.58}{m}$.
        So the ball is \final{\qty{4.6}{m} high} at the posts, above \qty{3.0}{m}:
        \textbf{yes, it goes over}, by about \qty{1.6}{m}. (Still in the air:
        $t_\mathrm{flight} = 2 \times 11.47/9.8 = \qty{2.34}{s}$.)
  % ---- 1.3 Energy and air resistance
  \item Energy is conserved and the mass cancels:
        $v = \sqrt{u^2 + 2gh} = \sqrt{9.0^2 + 2 \times 9.8 \times 1.5} = \sqrt{81 + 29.4}
        = \sqrt{110.4} = \qty{10.51}{m.s^{-1}}$.
        Answer: \final{$v = \qty{11}{m.s^{-1}}$} (2 significant figures), for any direction of throw.
  \item At the top the ball moves horizontally at $v_x = u\cos\ang{40}$. So
        $\dfrac{K_\mathrm{top}}{K_\mathrm{launch}} = \dfrac{\frac{1}{2}m(u\cos\ang{40})^2}{\frac{1}{2}mu^2}
        = \cos^2\ang{40}$. ($\cos^2\ang{40}$ means $(\cos\ang{40})^2 = 0.766^2$.) This is
        $0.587$, about \final{\qty{59}{\percent}}.
        The vertical motion has stopped (its energy is now potential energy), but the horizontal
        motion has not.
  \item Any three: (1) drag reduces the horizontal velocity, so the \textbf{range is shorter};
        (2) drag adds to gravity on the way up, so the \textbf{maximum height is lower}; (3) the
        path is \textbf{not symmetrical}: the way down is steeper and takes longer than the way
        up; (4) the ball \textbf{lands slower and at a steeper angle}; (5) the best angle for
        range is \textbf{less than \ang{45}}.
  % ---- 1.4 Predicting motion (added by the coverage audit;
  \item Constant acceleration, up positive, $a_y = \qty{-9.8}{m.s^{-2}}$.
        (a) At the top $v_y = 0$: $t_\mathrm{top} = \dfrac{u_y}{g} = \dfrac{9.0}{9.8} = \qty{0.918}{s}$
        (\final{\qty{0.92}{s}}).
        (b) $h_\mathrm{max} = \dfrac{u_y^2}{2g} = \dfrac{9.0^2}{2 \times 9.8} = \dfrac{81}{19.6}
        = \qty{4.13}{m}$ (\final{\qty{4.1}{m}}).
        (c) $v_y = u_y + a_y t = 9.0 + (-9.8)(1.5) = 9.0 - 14.7 = \qty{-5.7}{m.s^{-1}}$:
        \final{\qty{5.7}{m.s^{-1}} downwards}. The minus sign means it is falling; it passed the
        top at \qty{0.92}{s}.
  \item (a) Vertically the ball goes up and comes back to the same height:
        $t = \dfrac{2u_y}{g} = \dfrac{2 \times 3.5}{9.8} = \qty{0.714}{s}$ (\final{\qty{0.71}{s}}).
        (b) She has a uniform velocity: $s = ut = 6.0 \times 0.714 = \qty{4.29}{m}$
        (\final{\qty{4.3}{m}}).
        (c) \textbf{In her hand.} When the ball leaves her hand it already moves forwards at
        \qty{6.0}{m.s^{-1}}, like her. There is no horizontal force on it (air resistance
        ignored), so it keeps this uniform velocity. In \qty{0.71}{s} the ball and the cyclist
        both move \qty{4.3}{m} forwards, so the ball comes down where her hand is.
  \item (a) Gradient $= \dfrac{\text{rise}}{\text{run}} = \dfrac{-14.7 - 14.7}{3.0 - 0}
        = \dfrac{-29.4}{3.0} = \qty{-9.8}{m.s^{-2}}$. The gradient is the acceleration:
        \final{\qty{9.8}{m.s^{-2}} downwards}, the same at every moment.
        (b) The line crosses the time axis ($v = 0$, the top) half way, at \qty{1.5}{s}. The
        area of the triangle above the axis is
        $\tfrac12 \times \qty{1.5}{s} \times \qty{14.7}{m.s^{-1}} = \qty{11.0}{m}$, so the
        maximum height is \final{\qty{11}{m}}. (Check: $14.7^2 \div 19.6 = \qty{11.0}{m}$.)
        From \qtyrange{1.5}{3.0}{s} the triangle is below the axis, with the same area, so it
        counts as $-\qty{11.0}{m}$. Total: $11.0 + (-11.0) = 0$. The displacement at
        \qty{3.0}{s} is \final{zero}: the ball is back at its launch height.
  \item Marble, vertically (constant acceleration, $u_y = 0$):
        $-1.25 = \tfrac{1}{2}(-9.8)t^2 = -4.9t^2$. Divide both sides by $-4.9$:
        $t^2 = \dfrac{-1.25}{-4.9} = \qty{0.2551}{s^2}$, and
        $t = \qty{0.5051}{s}$.
        Marble, horizontally (uniform velocity): it lands $2.0 \times 0.5051 = \qty{1.010}{m}$
        from the bench.
        Cart (uniform velocity, towards the bench): in \qty{0.5051}{s} it moves
        $0.50 \times 0.5051 = \qty{0.2525}{m}$ closer.
        So when the marble leaves, the cart must be $1.010 + 0.2525 = \qty{1.263}{m}$ from the
        bench: \final{about \qty{1.3}{m}} (2 significant figures).
  % ---- 1.5 Key experiment
  \item Independent variable: the \textbf{launch angle $\theta$}. Dependent variable: the
        \textbf{range $R$}. Controlled variables (any two): the launch speed (the same spring
        setting); the same ball; launch and landing at the same height; the same landing
        surface; the same way of releasing the ball.
  \item The energy method assumes that all the lost potential energy becomes kinetic energy of
        the ball moving forward. But a rolling ball also spins, so some energy becomes kinetic
        energy of spinning (and a little is lost to rolling resistance and sound). So the real
        $u$ is smaller than calculated (about \qty{15}{\percent} smaller for a solid rolling
        ball), and the predicted ranges are too big. The experiment would then test the energy
        assumption, not the projectile model. A light gate measures $u$ directly, which makes
        it valid.
  \item A \textbf{systematic error}: every result is off in the same direction by about the
        same amount (random errors scatter both ways). Possible causes (any one): (1) the
        light-gate beam passes above or below the centre of the ball. The ball then blocks the
        beam for less than its full diameter, so the calculated $u$ is too big. (2) $h$ is
        measured too big. (3) The end of the ramp tilts slightly down. (Air resistance on a small steel ball at about \qty{2}{m.s^{-1}}
        changes $R$ by only about \qty{0.1}{\percent}, so it is not the main cause.)
\end{chapteranswers}
